%% Overleaf			
%% Software Manual and Technical Document Template	
%% 									
%% This provides an example of a software manual created in Overleaf.

\documentclass{../ol-softwaremanual}

% Packages used in this example
\usepackage{graphicx}  % for including images
\usepackage{microtype} % for typographical enhancements
\usepackage{minted}    % for code listings
\usepackage{amsmath}   % for equations and mathematics
\setminted{style=friendly,fontsize=\small}
\renewcommand{\listoflistingscaption}{List of Code Listings}
\usepackage{hyperref}  % for hyperlinks
\usepackage[a4paper,top=4.2cm,bottom=4.2cm,left=3.5cm,right=3.5cm]{geometry} % for setting page size and margins

\usepackage[english, greek]{babel}

\usepackage{subfig}

\usepackage{incgraph,tikz}

\usepackage{filemod}





\usepackage{rotating}


% Custom macros used in this example document
\newcommand{\doclink}[2]{\href{#1}{#2}\footnote{\url{#1}}}
\newcommand{\cs}[1]{\texttt{\textbackslash #1}}

\begin{document}
	
	
	\begin{titlepage}
		
		
		% Frontmatter data; appears on title page
		\title{\en Robustness Diagrams \\}
		\version{0.1}
		\softwarelogo{\includegraphics[scale=0.4]{../CarBazaar_logo.png}}		
		
	\end{titlepage}
	
	
	\maketitle
	
	\newpage
	
	\center{\textbf{Μέλη Ομάδας}}
	
	\vspace{20pt}
	
	
	
	\begin{table}[htbp!]
		
		\begin{tabular}{llll}
			Μεμελετζόγλου Χαρίλαος & 1069364 & \en st1069364@ceid.upatras.gr & 4o Έτος   \\ 
			\\ Λέκκας Γεώργιος      &      1067430    &   \en st1067430@ceid.upatras.gr & 4o Έτος  \\
			\\ Γιαννουλάκης Ανδρέας        &   1067387       & \en st1067387@ceid.upatras.gr & 4o Έτος           \\
			\\ Κανελλόπουλος Ιωακείμ        &  1070914        &    \en st1070914@ceid.upatras.gr & 4o Έτος        \\ 
		\end{tabular}
	\end{table}
	
	\center{\textbf{Υπεύθυνοι Παρόντος Τεχνικού Κειμένου}}
	
	\vspace{20pt}
	
	\begin{table}[htbp!]
		\begin{tabular}{ll}
			Μεμελετζόγλου Χαρίλαος & \en Editor \\
			\\ Λέκκας Γεώργιος      &   \en  Editor \\
			\\ Γιαννουλάκης Ανδρέας & \en Contributor \\
			\\ Κανελλόπουλος Ιωακείμ & \en Contributor \\ 
		\end{tabular}
	\end{table}
	
	
	\vspace{20pt}
	
	\center{\textbf{Εργαλεία που χρησιμοποιήθηκαν}}
	
	\vspace{20pt}
	\flushleft
	Χρησιμοποιήθηκε το \en \doclink{https://www.overleaf.com/}{Overleaf} \gr και το \en \doclink{https://www.texstudio.org/}{TexStudio} \gr για την συγγραφή του \LaTeX\ κώδικα. \break
	
	Για την δημιουργία του λογότυπου, χρησιμοποιήθηκε το εργαλείο \en \doclink{https://www.adobe.com/express/create/logo}{Adobe Express} . \gr \break
	
	Για την δημιουργία των \en Robustness Diagrams \gr χρησιμοποιήθηκε το \en \doclink{https://www.visual-paradigm.com/}{Visual Paradigm} . \gr \break 
	
	\newpage
	
	\center{\textbf{\en Robustness Diagrams \gr}}
	\flushleft
	
	Οι εναλλακτικές ροές του κάθε \en Use Case \gr φαίνονται στο αντίστοιχο \en Robustness Diagram \gr, με κόκκινες ακμές και αντικείμενα. \break
	
	Για την ευκολότερη αντιστοίχιση ενός \en Use Case \gr στο αντίστοιχο \en Robustness Diagram \gr, παραθέτουμε και το κείμενο της κάθε Περίπτωσης Χρήσης, πρίν το διάγραμμα που προκύπτει. \\
	
	\newpage
	
	\centering
	\paragraph{\en Use Case 2: \gr Προγραμματισμός Ελέγχου Οχήματος} 
	
	\begin{enumerate}
		\item Ο χρήστης επιλέγει \en"\gr Έλεγχος Οχήματος\en" \gr στο αρχικό μενού
		\item Το σύστημα εμφανίζει την σελίδα Προγραμματισμού Ελέγχου Οχήματος
		\item Ο χρήστης επιλέγει το πακέτο ελέγχου που επιθυμεί, την ημερομηνία και ώρα διεξαγωγής του ελέγχου και εισάγει την τοποθεσία του
		\item Το σύστημα αφού επιβεβαιώσει την εισαχθείσα τοποθεσία, προτείνει στην χρήστη έναν ελεγκτή, με βάση την τοποθεσία που εισήγαγε
		\item Ο χρήστης αποδέχεται ή όχι τον προτεινόμενο ελεγκτή
		\item Το σύστημα εμφανίζει την οθόνη Εισαγωγή Κωδικού Αγγελίας, στο όχημα της οποίας θα πραγματοποιηθεί ο έλεγχος		
		\item Ο χρήστης εισάγει τον κωδικό της αγγελίας
		\item Το σύστημα ανακτά τα στοιχεία του οχήματος από την αγγελία και εμφανίζει την τελική τιμή του ελέγχου καθώς και την χρονική διάρκειά του
		\item Ο χρήστης επιβεβαιώνει τα στοιχεία
		\item Το σύστημα μεταφέρει τον χρήστη στο μενού πληρωμών. Μετά την επιτυχή πληρωμή, εμφανίζει μήνυμα επιτυχούς κράτησης και αποστέλλει \en email \gr στον χρήστη, με τα στοιχεία του ραντεβού και του ελεγκτή		
	\end{enumerate}
	
	\paragraph{Εναλλακτική Ροή 1}
	
	\begin{enumerate}
		\item Ο χρήστης εισάγει μη-υπαρκτή τοποθεσία
		\item Το σύστημα ενημερώνει τον χρήστη με το κατάλληλο μήνυμα σφάλματος 
		\item Ο χρήστης εισάγει την σωστή τοποθεσία του
		\item Το σύστημα εντοπίζει τον χρήστη και η Περίπτωση Χρήσης προχωρά από το βήμα 4 της βασικής ροής
	\end{enumerate}
	
	\paragraph{Εναλλακτική Ροή 2}
	
	\begin{enumerate}
		\item Ο χρήστης απορρίπτει τον προτεινόμενο από το σύστημα ελεγκτή, προκειμένου να επιλέξει τον ελεγκτή της αρεσκείας του
		\item Το σύστημα εμφανίζει την οθόνη εισαγωγής στοιχείων του ελεγκτή
		\item Ο χρήστης εισάγει τα στοιχεία του ελεγκτή
		\item Το σύστημα ελέγχει πως υπάρχει πράγματι εγγεγραμμένος ο εν λόγω ελεγκτής και αν ναι, η Περίπτωση Χρήσης συνεχίζει από το βήμα 6 της βασικής ροής. Ειδάλλως, εμφανίζει μήνυμα σφάλματος και ο έλεγχος ακυρώνεται	
	\end{enumerate}	


	\begin{figure}[htbp!]
		\includegraphics[scale=0.45]{img/rob\_car\_check.png}
		\caption{\en Robustness Diagram : "\gr Προγραμματισμός Ελέγχου Οχήματος\en"\gr}
	\end{figure}
	
	
	
	
	
\end{document}



 
